\begin{table}

\caption{Inclusion of Hurricane Maria. \label{tab:appB_maria}}
\centering
\begin{adjustbox}{max width = \textwidth}
\begin{threeparttable}
\begin{tabular}[t]{lrrrr}
\toprule
Sample & Obs. & $\delta$ (m/s) & SE & p\\
\midrule
Excluding Maria (baseline) & 124 & 0.258 & (0.095) & 0.007\\
Including Maria & 274 & 0.356 & (0.045) & $<$0.001\\
\bottomrule
\end{tabular}
\begin{tablenotes}
\small
\item Both rows estimate the Stage 2 precision-gain specification of Table \ref{tab:noaa-controls} on hurricane-force ($\geq$ 33 m/s) realizations at the 3-day lead. The baseline row uses the six 2017--2020 continental-US landfalls in the AO denial archive; the second row adds Hurricane Maria back into the panel. The coefficient $\delta$ is the average reduction in the ex-ante wind-speed forecast standard deviation attributable to the assimilation of dropsonde profiles (the ALL configuration of \citet{DitchekEA2023}). Maria is excluded from the main analysis because it did not make landfall in the continental US. County-clustered standard errors in parentheses.
\end{tablenotes}
\end{threeparttable}
\end{adjustbox}
\end{table}
